\section{Introduction}

\subsection{Background}
It was noticed a little more than a decade ago \cite{EM13, FPSS12} that two seminal theorems of Erd\H{o}s and Szekeres from their 1935 paper \cite{ES35}---the \textit{monotone subsequence theorem} and the \textit{cups-caps lemma}---can be cast as Ramsey theoretic statements about monotone paths in two- and three-uniform hypergraphs, respectively.
A \textit{tight monotone path of length $n$} in a $k$-uniform hypergraph on $[N]$ is the set of $n$ edges which contain $k$ consecutive vertices from $v_1 < v_2 < \dots < v_{n+k-1}$ in $[N]$, and we will denote by $MS_k(n; r)$ the least integer $N$ such that every $r$-coloring of $[N]^{(k)}$ contains a tight \textit{monochromatic} monotone path of length $n$, i.e. a monotone path in which every edge receives the same color.
(Here and throughout, for a natural number $N$, we will denote by $[N]$ the set $\{1,2,\dots,N\}$, and for a set $X$ and a natural number $k$, by $X^{(k)}$ the set $\{e \subseteq X: |e| = k\}$.)
In this language, the theorems of Erd\H{o}s and Szekeres read as follows.

\begin{theorem}[Erd\H{o}s--Szekeres, 1935 \cite{ES35}]
    For all positive integers $n \geq 3$ and $r \geq 2$, we have
    \begin{equation*}
        MS_2(n; 2) = n^2 + 1.
    \end{equation*}
    Also, for $n \geq 4$,
    \begin{equation*}
        MS_3(n; 2) = \binom{2n}{n} + 1.
    \end{equation*}
\end{theorem}

In \cite{FPSS12}, the authors identified this unifying framework and asked for sharp estimates on $MS_k(n; r)$ for larger values of $k$ and $r$.
Soon after, the problem was almost entirely resolved by Moshkovitz and Shapira \cite{MS14}, who established a connection between Ramsey numbers of monotone paths and certain types of higher-order down-sets as well as high-dimensional integer partitions.
They proved an exact equality between $MS_k(n; r)$ and the number of these down sets which yielded the following quantitative estimates.
To state them, we recursively define the \textit{tower function} $T_k(m)$ by $T_0(m) = m$ and $T_{k+1}(m) = 2^{T_k(m)}$; note that the subscript denotes the number of 2's in the tower.

\begin{theorem}[Moshkovitz--Shapira, 2014 \cite{MS14}]\label{MS}
    For all $k \geq 3$, $r \geq 2$, and $n \geq 2$,
    \begin{equation*}
        MS_k(n; r) \leq T_{k-2}(2 n^{r-1}).
    \end{equation*}
    Moreover, for $n$ sufficiently large,
    \begin{equation*}
        MS_k(n; r) \geq T_{k-2}(n^{r-1} / 2 \sqrt{r}).
    \end{equation*}
\end{theorem}

